\documentclass[conference]{IEEEtran}
\usepackage[T1]{fontenc}
\usepackage{amsmath,amssymb,amsfonts}
\usepackage{bm}
\usepackage{booktabs}
\usepackage{graphicx}
\usepackage{cite}
\usepackage{url}
\usepackage{algorithm}
\usepackage{algorithmic}

\newcommand{\KVTTT}{\textsc{KV-TTT}}

\begin{document}

\title{When Can Internal-State Adaptation Replace Weight Adaptation in Vision--Language Models?}
\author{\IEEEauthorblockN{Anonymous Authors}}
\maketitle

\begin{abstract}
Few-shot adaptation of a vision--language model (VLM) commonly changes model
weights, even when the target shift may require only a small correction to how
visual evidence is read.  We ask a broader question: \emph{when can adapting a
frozen model's internal state replace weight adaptation?}  We study a simple
test-time method, \emph{Key--Value Test-Time Tuning} (\KVTTT{}), that constructs
a low-rank basis from the second moment of support correctness gradients and
learns a bounded residual on task-relevant visual key/value activations.  By
the Ky Fan maximum principle, this basis is the rank-$r$ subspace retaining the
maximum support gradient energy.  Across disaster assessment, histopathology,
robot failure recognition, and vision--language action classification, we find
that energy coverage alone is not predictive.  Successful internal-state
adaptation additionally requires signed support--query directional transfer
and an activation site with sufficient functional authority.  On fixed-
semantics \textsc{BRIGHT}, \KVTTT{} exceeds matched LoRA on average with
$1{,}024$ learned coefficients versus $10.1$M LoRA parameters, and the pattern
replicates across VLM families.  Camelyon exhibits stable geometry and a
query-basis KV oracle that reaches LoRA, while RoboFail exposes class-mixture
and site mismatch: oracle prior reweighting raises directional agreement from
$0.209\!\pm\!0.361$ to $0.934\!\pm\!0.025$.  ManipBench remains far below LoRA
even with query-derived bases and broader activation sites, identifying a
capability boundary rather than an optimization failure.  These results turn
internal-state adaptation from a universal replacement claim into a testable
account of when local, calibration-preserving adaptation is sufficient.
\end{abstract}

\begin{IEEEkeywords}
test-time adaptation, vision--language models, internal-state adaptation,
gradient geometry, parameter-efficient adaptation.
\end{IEEEkeywords}

\section{Introduction}
\IEEEPARstart{F}{ew-shot} deployment poses an ambiguous adaptation problem.  A
new domain may preserve a capability already present in a VLM and change only
which visual evidence should be trusted.  Alternatively, it may require a new
language--vision decision rule or action mapping.  Weight-space methods such as
LoRA~\cite{hu2021lora} can address both cases, but use far more degrees of
freedom than necessary in the first and can produce severely overconfident
errors.  Internal-state adaptation is a compelling smaller alternative, yet
its applicability conditions are poorly understood.

We ask whether event adaptation can instead be expressed as a low-dimensional
change in a frozen model's working activations. Transformer key/value projections
form a natural intervention point~\cite{vaswani2017attention,geva2021kv}. For
each labeled support example, we differentiate the answer-token loss with
respect to the post-event image-token projection outputs. Crucially, we retain
the second moment $\sum_iG_i^\top G_i$, not the average gradient. Its leading
eigenvectors preserve a space of corrective directions even when example-level
gradients cancel in the mean.

Our central hypothesis is that successful internal-state test-time adaptation
requires two conditions: \emph{directional geometry transfer}, meaning that
support and query request compatible corrections in the learned subspace, and
\emph{actuator sufficiency}, meaning that the intervened activation can realize
the target functional change.  We test this hypothesis rather than expanding
the method until it approximates a small weight adapter.

Our contributions are: (i) a frozen-backbone, visual-KV residual whose
gradient-covariance basis is optimal for rank-constrained support energy
retention; (ii) energy-overlap ($\rho$), signed-alignment ($\kappa$),
class-conditional, query-basis, and projection-site diagnostics that separate
geometry from actuation; (iii) evidence across \textsc{BRIGHT}, Camelyon17,
Guardian/FailCoT, and ManipBench, including multiple VLM families and matched
three-seed task evaluations; and (iv) explicit positive, mismatch, and
capability-boundary regimes.  Query labels are used only in clearly marked
mechanism oracles, never for primary adaptation or model selection.

\section{Related Work}
\textbf{Test-time adaptation.} Test-time training~\cite{sun2020ttt},
TENT~\cite{wang2020tent}, CoTTA~\cite{wang2022cotta}, and test-time prompt
tuning~\cite{shu2022tpt} adapt normalization, objectives, or prompts at
deployment. \KVTTT{} instead uses a small labeled event support once, produces a
fixed event controller, and adds no query context.

\textbf{Parameter-efficient and activation-space adaptation.} LoRA and
adapters modify weights~\cite{hu2021lora,houlsby2019adapters}; prompt tuning
learns input-like parameters~\cite{lester2021prompt}. Transformer layers have
also been interpreted and edited as key--value memories
~\cite{geva2021kv,meng2022rome,meng2023memit}. Our method neither edits nor
adds model weights: it applies a bounded residual to selected activation rows.

\textbf{Remote-sensing damage assessment.} xView2~\cite{gupta2021xview} and
\textsc{BRIGHT}~\cite{bright2024} support building-level damage research, while
remote-sensing foundation representations and VLMs broaden transfer
~\cite{ravi2023satmae,bhatia2024rerank}. We focus on same-event, few-example
adaptation of a general VLM under strict spatial separation.

\section{Method}
\subsection{Gradient-covariance subspace}
Let $f_\theta$ be a frozen Qwen2.5-VL-7B model~\cite{qwen2_5vl}. For support
example $(x_i,y_i)$, projection type $t\in\{K,V\}$, and selected decoder layer
$l$, let $Z_i^{(l,t)}\in\mathbb{R}^{T\times d}$ denote the projection output.
A binary mask $M_i$ selects only the second (post-event) image-token group. We
compute the answer-span loss and masked gradient
\begin{equation}
G_i^{(l,t)}=\nabla_{[Z_i^{(l,t)}]_{M_i}}\mathcal{L}(x_i,y_i).
\end{equation}
The uncentered gradient second moment and rank-$r$ basis are
\begin{equation}
C^{(l,t)}=\frac{1}{N}\sum_{i=1}^{N}(G_i^{(l,t)})^\top G_i^{(l,t)},\quad
B^{(l,t)}=\operatorname{TopEig}_r(C^{(l,t)}).
\end{equation}
This is equivalent to retaining the leading right singular directions of the
stacked token gradients. It preserves directional variation that
$\bar g=N^{-1}\sum_i g_i$ discards.

\paragraph{Optimal rank-constrained energy retention.}
Let $C\succeq0$ and let $B\in\mathbb{R}^{d\times r}$ satisfy
$B^\top B=I_r$.  The correctness-gradient energy retained by the subspace is
\begin{equation}
\mathcal E(B)=\operatorname{tr}(B^\top C B)
=\frac{1}{N}\sum_i\lVert G_iB\rVert_F^2.
\end{equation}
By the Ky Fan maximum principle~\cite{fan1951maximum},
\begin{equation}
B^*=\operatorname{TopEig}_r(C)
\in\arg\max_{B^\top B=I_r}\operatorname{tr}(B^\top C B),
\end{equation}
and the optimum equals the sum of the largest $r$ eigenvalues of $C$.
Thus the basis is not an arbitrary SVD heuristic: among all rank-$r$
subspaces it preserves the most support correctness-gradient energy.  This
property is deliberately limited.  It guarantees neither that query gradients
occupy the same subspace, that their signed mean agrees with support, nor that
a residual at the chosen activation can express the required decision change.

\paragraph{Geometry and actuator diagnostics.}
For support basis $B_S$ and stacked query gradients $G_Q$, we measure energy
transfer and signed agreement as
\begin{equation}
\rho=\frac{\lVert G_QB_S\rVert_F^2}{\lVert G_Q\rVert_F^2},\qquad
\kappa=\cos(P_{B_S}\bar g_S,P_{B_S}\bar g_Q),
\end{equation}
where $P_{B_S}=B_SB_S^\top$.  High $\rho$ says that relevant query energy is
covered; high $\kappa$ says that the requested average correction is aligned.
We then hold subspace size fixed and vary the intervention site or use a
query-derived basis as an oracle test of actuator authority.  These diagnostics
read query labels and are never deployable performance estimates.

\subsection{Bounded residual controller}
For a Diagonal controller,
\begin{equation}
Z'=Z+M\odot\left[(ZB)\operatorname{diag}(\alpha\tanh a)B^\top\right].
\end{equation}
For a Full controller, a raw $r\times r$ matrix $R$ gives
\begin{equation}
A=\frac{\alpha R}{1+\lVert R\rVert_F},\qquad
Z'=Z+M\odot\left[(ZB)AB^\top\right].
\end{equation}
Thus $\lVert A\rVert_2<\alpha$. All coefficients start at zero, so the
controller initially implements the identity. We use the middle and final
decoder layers (14 and 27); both $K$ and $V$ are controlled. Rank-16 Full and
Diagonal variants learn $2\times2\times16^2=1{,}024$ and
$2\times2\times16=64$ scalars, respectively. The compact rank-5 variants learn
100 and 20.

\subsection{Coefficient fitting}
The frozen VLM and bases remain fixed while Adam optimizes only controller
coefficients on all 24 support examples:
\begin{equation}
\min_\phi\frac{1}{N}\sum_i\mathcal{L}(f_{\theta,B,\phi}(x_i),y_i)
+\lambda\lVert\phi\rVert_2^2.
\end{equation}
The common configuration uses four full-support updates, learning rate $0.05$,
$\lambda=10^{-3}$, and unit gradient clipping. Each query then runs normally
with the fixed controller; no query-time optimization is performed.

\begin{algorithm}[t]
\small
\caption{Supervised event-level \KVTTT{}}
\label{alg:kvttt}
\begin{algorithmic}[1]
\REQUIRE Frozen VLM $f_\theta$; support $\mathcal S$; layers $\mathcal L$;
rank $r$; bound $\alpha$; updates $T_u$
\FOR{$(x_i,y_i)\in\mathcal S$}
  \STATE Build the paired-image prompt and isolate the answer-token span.
  \STATE Backpropagate its loss; retain post-image gradients $G_i^{(l,t)}$.
  \STATE $C^{(l,t)}\leftarrow C^{(l,t)}+(G_i^{(l,t)})^\top G_i^{(l,t)}$.
\ENDFOR
\FOR{$l\in\mathcal L$ and $t\in\{K,V\}$}
  \STATE $B^{(l,t)}\leftarrow\operatorname{TopEig}_r(C^{(l,t)})$.
  \STATE Initialize the Diagonal vector $a$ or Full matrix $R$ to zero.
\ENDFOR
\FOR{$u=1,\ldots,T_u$}
  \STATE Accumulate the mean labeled loss over all $|\mathcal S|=24$ items.
  \STATE Add $\lambda\lVert\phi\rVert_2^2$, clip gradients, and update $\phi$.
\ENDFOR
\RETURN Bases $B$ and controller coefficients $\phi$.
\end{algorithmic}
\end{algorithm}

\paragraph{Complexity.}
For $L$ selected layers and both projection types, Diagonal learns $2Lr$
scalars and Full learns $2Lr^2$. Basis extraction requires one frozen-model
backward pass per support item and an eigendecomposition of each $d\times d$
second-moment matrix. The controller adds two low-rank projections per selected
module at inference but no extra tokens, retrieval, or optimization. Bases and
coefficients are serialized together; their storage is dominated by $B$, not
by the learned coefficient count.

\section{Experiments}
\subsection{Protocol}
We use the four \textsc{BRIGHT} events for which the strict balanced protocol
is complete: Hawaii wildfire, Libya flood, Noto earthquake, and Turkey
earthquake. Each event has 24 same-event labeled support instances (8 per
class) and 300 query instances (100 per class); support and query tile sets are
disjoint. All arms start independently from raw
\texttt{Qwen/Qwen2.5-VL-7B-Instruct}, use $448\times448$ paired crops, and score
the same query IDs. The clean LoRA baseline uses rank 16, alpha 32,
$q/k/v/o$ targets, learning rate $10^{-4}$, batch size 1, gradient accumulation
3, and eight optimizer updates---exactly one pass over support24, without
support cross-validation. Metrics are macro-F1, balanced accuracy, and NLL.

\paragraph{Cross-task evaluation.}
We additionally evaluate InternVL3-8B and Qwen3-VL-8B on Camelyon17 hospital
2 (binary histopathology, support16/query300) and ManipBench Q1 (four-way
vision--language action classification, query400), with three matched support
seeds and identical query IDs for Frozen, rank-16 LoRA, Random-KV, and
\KVTTT{}. Guardian/FailCoT uses InternVL3-8B on three fixed-semantics
success/failure splits, support16, and three matched seeds. Corrected
InternVL3 task results use the required \texttt{Answer: <label>} format; older
raw-label ManipBench runs are excluded. Main performance results never use
query labels. Mechanism diagnostics do, and are reported separately as
oracles.

The evidence blocks serve different statistical roles. BRIGHT's primary
four-method comparison uses one locked support24 per event; its independent
three-support-seed suite estimates KV-versus-Frozen robustness but does not
provide a matched multiseed LoRA comparison. Camelyon, Guardian, and
ManipBench use matched support seeds. We keep these estimands separate rather
than pooling incompatible replications.

\subsection{Fixed-configuration main comparison}
Table~\ref{tab:fixed} uses one common rank-16, alpha-3, four-update KV
configuration across events. Full improves the raw VLM on every event and has
the best mean F1, while LoRA has the best mean NLL. Full exceeds fixed-8 LoRA
on Libya and Turkey; LoRA remains better on Hawaii and Noto. Diagonal uses only
64 learned scalars but is less accurate. The table exposes the intended
accuracy--adaptation-size trade-off rather than hiding metrics on which LoRA is
better.

\begin{table*}[t]
\centering\small
\caption{Fixed-configuration comparison. Event columns and their mean report
Macro-F1; NLL is averaged over events (lower is better). KV uses rank 16,
layers 14+27, alpha 3, and four updates. ``Learned'' counts event-specific
optimized scalars, excluding frozen covariance bases.}
\label{tab:fixed}
\begin{tabular}{lrrrrrrr}
\toprule
Method & Hawaii & Libya & Noto & Turkey & Mean F1 $\uparrow$ & Mean NLL $\downarrow$ & Learned $\downarrow$\\
\midrule
Raw VLM & 0.1992 & 0.2585 & 0.2087 & 0.2273 & 0.2234 & 1.3233 & 0\\
LoRA fixed8 & \textbf{0.3526} & 0.2776 & \textbf{0.3495} & 0.2147 & 0.2986 & \textbf{1.1276} & 10,092,544\\
Diagonal KV & 0.2918 & 0.2560 & 0.2820 & 0.2879 & 0.2794 & 1.1715 & \textbf{64}\\
Full KV & 0.2957 & \textbf{0.3739} & 0.3112 & \textbf{0.3025} & \textbf{0.3208} & 1.1367 & 1,024\\
\bottomrule
\end{tabular}
\end{table*}


\begin{table}[t]
\centering\small
\caption{Four-event BRIGHT mean macro-F1 under the locked support24/query300
protocol for two additional 8B backbones.}
\label{tab:brightbackbones}
\begin{tabular}{lrrrr}
\toprule
Backbone & Frozen & LoRA & Random-KV & Ours\\
\midrule
Qwen3-VL-8B & .1667 & .2034 & .1687 & \textbf{.3492}\\
InternVL3-8B & .1848 & .1667 & .2214 & \textbf{.2860}\\
\bottomrule
\end{tabular}
\end{table}

These architecture-transfer results reproduce the BRIGHT ranking without
selecting a different actuator.  Additional Phi, Gemma, and LLaVA controls are
mixed---Ours is strongest on Gemma, LoRA on Phi, and Frozen on LLaVA---so the
claim is cross-architecture evidence, not universal dominance.

\begin{table*}[t]
\centering\small
\caption{Corrected InternVL3 cross-task macro-F1 (mean $\pm$ sample SD over
three matched support seeds). Guardian reports the fixed eight-step Ours
follow-up; Random-KV was not run in that suite.}
\label{tab:crosstask}
\begin{tabular}{lrrrr}
\toprule
Task & Frozen & LoRA & Random-KV & Ours\\
\midrule
Camelyon17 hospital 2 & $.3333\pm.0000$ & $.6325\pm.2725$ & $.4935\pm.0071$ & $.5107\pm.0837$\\
Guardian UR5-Fail & $.4513\pm.0067$ & $.4432\pm.0510$ & -- & $\mathbf{.4905\pm.0308}$\\
Guardian RoboFail & $.4476\pm.0263$ & $\mathbf{.5010\pm.0645}$ & -- & $.3952\pm.0310$\\
Guardian RoboVQA & $.4043\pm.0068$ & $.5159\pm.0087$ & -- & $\mathbf{.5204\pm.0108}$\\
ManipBench droid-pick & $.1541\pm.0000$ & $\mathbf{.4304\pm.2223}$ & $.1556\pm.0047$ & $.2839\pm.0729$\\
\bottomrule
\end{tabular}
\end{table*}

\subsection{Rank-16 alpha sweep}
We sweep alpha in $\{0.5,1,2,3,5,10\}$ for both controllers on all events
(48/48 runs). Table~\ref{tab:oracle} reports per-event maxima to expose method
capacity. These are \emph{query-oracle diagnostics}, not valid primary results:
alpha must be chosen using support-only evidence before query evaluation.

\begin{table}[t]
\centering\small
\caption{Per-event query-oracle rank-16 alpha maxima.}
\label{tab:oracle}
\begin{tabular}{lrrrr}
\toprule
Event & LoRA & Full & $\alpha$ & Diag. / $\alpha$\\
\midrule
Hawaii & 0.3526 & \textbf{0.3869} & 2 & 0.3341 / 10\\
Libya & 0.2776 & \textbf{0.3852} & 2 & 0.3390 / 5\\
Noto & \textbf{0.3495} & 0.3270 & 1 & 0.3404 / 5\\
Turkey & 0.2147 & \textbf{0.3025} & 3 & 0.2879 / 3\\
\midrule
Mean & 0.2986 & \textbf{0.3504} & -- & 0.3254 / --\\
\bottomrule
\end{tabular}
\end{table}

An independent-repeat warning is important: two nominally identical Hawaii
rank-16 Full alpha-3 runs produce macro-F1 0.3295 and 0.2957 despite matching
model fingerprint and recorded configuration. We therefore do not promote the
oracle alpha-2 peak to the primary claim until repeated.

\subsection{Mechanism controls and statistical robustness}
At rank 5 and alpha 3, uncentered covariance reaches mean macro-F1 0.2964,
versus 0.2506 for a rank-1 mean gradient, 0.2721 for centered covariance, and
0.2414 for a three-seed random basis. Ten-thousand-resample paired tests give
covariance-minus-mean $+0.0458$ ($p=0.0142$), covariance-minus-centered
$+0.0243$ ($p=0.0032$), and covariance-minus-random $+0.0551$ ($p=0.0022$).
One random seed approaches covariance, so the evidence supports higher and
more reliable average performance, not failure of every random subspace.

Across three balanced support seeds and four events, compact Full versus raw
VLM improves macro-F1 by $+0.0457$ (95\% CI $[+0.0046,+0.0847]$,
$p=0.0167$); Diagonal improves by $+0.0399$ (95\% CI
$[+0.0063,+0.0734]$, $p=0.0061$). Full versus Diagonal is not significant
($p=0.7347$).

Table~\ref{tab:significance} places both performance and mechanism tests in the
main paper. Confidence intervals and permutation tests operate on paired event
and query outcomes; the random-basis comparison additionally samples the random
basis seed within each bootstrap replicate. All reported improvements over raw
and all covariance-mechanism contrasts reject a zero mean difference at the
5\% level, whereas Full and Diagonal are statistically indistinguishable.

\begin{table*}[t]
\centering\small
\caption{Paired statistical validation (10,000 resamples). Positive $\Delta$
favors the first method. Mechanism rows use the fixed seed-0 rank-5 controller;
raw comparisons sample three support seeds.}
\label{tab:significance}
\begin{tabular}{lrrr}
\toprule
Comparison & Mean $\Delta$ Macro-F1 & 95\% bootstrap CI & Permutation $p$\\
\midrule
Full KV $-$ raw VLM & $+0.0457$ & $[+0.0046,+0.0847]$ & $0.0167$\\
Diagonal KV $-$ raw VLM & $+0.0399$ & $[+0.0063,+0.0734]$ & $0.0061$\\
Full KV $-$ Diagonal KV & $+0.0058$ & $[-0.0433,+0.0506]$ & $0.7347$\\
Covariance $-$ mean gradient & $+0.0458$ & $[+0.0188,+0.0735]$ & $0.0142$\\
Covariance $-$ centered covariance & $+0.0243$ & $[+0.0100,+0.0388]$ & $0.0032$\\
Covariance $-$ random basis & $+0.0551$ & $[+0.0058,+0.1036]$ & $0.0022$\\
\bottomrule
\end{tabular}
\end{table*}

\subsection{When does internal-state adaptation suffice?}
Table~\ref{tab:regimes} combines primary performance with oracle mechanism
evidence.  It is not a correlation fit: the rows have deliberately different
roles, and BRIGHT geometry uses archived support12 folds while its primary
performance uses support24.  Instead, the table implements a sequence of
falsification tests for geometry transfer and actuator sufficiency.

\begin{table*}[t]
\centering\small
\caption{Geometry--Actuator regimes.  $\Delta$F1 is Ours minus LoRA under the
corresponding clean primary protocol.  Query-basis and site results are
query-label oracles, not deployable scores.}
\label{tab:regimes}
\begin{tabular}{lrrrrl}
\toprule
Task & $\rho$ & $\kappa$ & $\Delta$F1 & Same-actuator oracle & Diagnosis\\
\midrule
BRIGHT Hawaii & .433 & .975 & +.171 & -- & fixed-semantics evidence shift\\
BRIGHT Libya & .369 & .677 & +.058 & -- & transferable class geometry\\
Camelyon17 & $.785\pm.005$ & $.977\pm.008$ & -.122 & query-basis KV .663 vs LoRA .633 & finite-support/margin gap\\
RoboFail & $.620\pm.008$ & $.209\pm.361$ & -.106 & Q-only .481 vs LoRA .501 & mixture and site mismatch\\
ManipBench droid-pick & $.620\pm.003$ & $.741\pm.142$ & -.147 & query-basis KV .248 vs LoRA .430 & actuator/capability boundary\\
\bottomrule
\end{tabular}
\end{table*}

\paragraph{Positive evidence-shift regime.}
BRIGHT keeps the damage labels and decision semantics fixed while geography,
sensor conditions, and damage appearance change.  On the four-event locked
Qwen2.5 protocol, Full KV obtains .3208 mean macro-F1 versus .2986 for LoRA and
.2234 frozen, using 1,024 learned coefficients rather than 10.1M weight
parameters.  Qwen3-VL and InternVL3 show the same qualitative ordering in the
four-method architecture-transfer table.  Hawaii has strongly positive
directional agreement in every class; Libya's aggregate agreement is weaker,
but its class kappas all exceed .94.

\paragraph{Geometry transfers but realization can fail.}
Camelyon has stable high $\rho$ and $\kappa$ across three seeds.  Replacing the
support basis by a query-derived basis while retaining the same KV actuator
reaches .663 F1, comparable to or above the LoRA mean.  The local actuator can
therefore express the correction, while finite-support subspace estimation and
coefficient-margin realization remain error sources.  This is also visible in
calibration: corrected InternVL \KVTTT{} has NLL .784 versus 1.990 for LoRA
despite lower macro-F1, consistent with conservative probability correction
rather than aggressive boundary movement.

\paragraph{Mixture and site mismatch.}
RoboFail's class-conditional kappas are stable across seeds (success .962,
failure .970), whereas aggregate kappa is seed-sensitive.  Keeping the
support covariance and basis fixed but reweighting support class means with
the oracle query prior changes aggregate kappa from $.209\pm.361$ to
$.934\pm.025$.  The recovery occurs for Q, K, V, and O, directly implicating
class-mixture cancellation.  Q-only is the strongest seed-0 site diagnostic,
but is not promoted to the method because it is not consistently beneficial
elsewhere.

\paragraph{Capability boundary.}
ManipBench asks the model to compose instructions and images into changing
A/B/C/D action semantics.  Its directional agreement is positive, yet a
query-derived KV basis reaches only .248 F1 and QKVO only .263, far below the
.430 LoRA mean.  Longer coefficient optimization improves ordinary KV but
does not change this ordering.  The bottleneck is therefore not merely finding
the direction: local activation control lacks the functional authority of
weight adaptation for task/action remapping.

\subsection{Capacity, support sensitivity, and efficiency}
Adapting every decoder layer with rank-16 Full matrices (14,336 scalars) is not
beneficial. Naively retaining alpha 3 and learning rate 0.05 causes support-loss
spikes and mean performance below the sparse controller. Reducing alpha to 0.3
and learning rate to 0.005 stabilizes Hawaii support loss, but query macro-F1
is only 0.2923 and damaged recall is zero. Sparse middle+final intervention is
therefore empirically preferable to dense all-layer control.

Support selection also matters. Hawaii's original balanced support24 spans
only three tiles. A representative-medoid alternative remains balanced and
query-tile-disjoint but obtains 0.3084, below the original-support rank-5
alpha-2 value 0.3278. We report this as sensitivity, not as a replacement split.

On one RTX 3090 and the same 300 Hawaii queries, raw VLM, compact Full, and
compact Diagonal require 1.3405, 1.3382, and 1.3384 seconds/sample; LoRA requires
1.3725 seconds/sample. Peak memory is 17,470 MiB for raw and KV and 17,508 MiB
for LoRA. Compact Full and Diagonal learn 100 and 20 scalars (serialized states
approximately 44 KiB), versus $10{,}092{,}544$ parameters and 39.5 MiB for
LoRA. KV extraction and coefficient fitting average 17 and 53 seconds.

\section{Limitations}
BRIGHT's clean four-method table uses one locked support set per event.  Its
three-seed robustness suite repeats KV but not LoRA, so we do not claim a
matched multiseed KV--LoRA significance test.  The archived BRIGHT geometry
folds use support12 rather than the primary support24 and provide qualitative,
not correlational, evidence.  Query-basis, query-prior, and projection-site
experiments use query labels and are mechanism oracles only.  Small supports
remain sensitive to spatial and class-mixture composition.  The evaluation
spans several VLM families, but not all backbones are run on every task.
Finally, local KV control is intentionally not universal: it remains weaker
than LoRA when the target requires substantial action/semantic remapping.
These boundaries are part of the claim rather than targets for dataset-
specific actuator selection.

\section{Conclusion}
Gradient-covariance \KVTTT{} shows that a frozen VLM can sometimes be adapted
with a tiny visual activation-state controller instead of a weight adapter.
The key qualification is mechanistic: maximum support gradient-energy
retention is only the first step.  The correction must also transfer with the
right sign, and the selected activation must have enough causal authority to
realize it.  BRIGHT illustrates the positive fixed-semantics evidence-shift
regime; Camelyon separates actuator capacity from finite-support realization;
RoboFail exposes class-mixture and site mismatch; and ManipBench establishes a
task-remapping boundary.  The resulting answer to ``when can internal-state
adaptation replace weight adaptation?'' is neither always nor never: when the
target correction is directionally transferable and locally expressible.

\appendices
\section{Complete Fixed-Configuration Metrics}
Table~\ref{tab:allmetrics} reports balanced accuracy (BA), NLL, and ordinal MAE
in addition to the macro-F1 values in Table~\ref{tab:fixed}. A lower NLL and MAE
is better. The results show that improvements in macro-F1 need not imply an
ordinal-error improvement: for example, fixed-8 LoRA raises Hawaii F1 while
increasing MAE. This is why we retain both classification and probabilistic
metrics rather than selecting a method from F1 alone.

\begin{table*}[t]
\centering\scriptsize
\caption{Complete metrics for the fixed rank-16, alpha-3 comparison.}
\label{tab:allmetrics}
\begin{tabular}{llrrrr}
\toprule
Event & Method & Macro-F1 & BA & NLL & Ordinal MAE\\
\midrule
Hawaii & Raw / LoRA / Full / Diag.
& 0.1992 / 0.3526 / 0.2957 / 0.2918
& 0.3433 / 0.3767 / 0.3533 / 0.3300
& 1.2491 / 1.0974 / 1.1314 / 1.1235
& 0.6767 / 0.8600 / 0.9467 / 0.7800\\
Libya & Raw / LoRA / Full / Diag.
& 0.2585 / 0.2776 / 0.3739 / 0.2560
& 0.3633 / 0.3467 / 0.3800 / 0.3567
& 1.3331 / 1.1175 / 1.1011 / 1.1525
& 0.6633 / 0.9200 / 0.8167 / 0.6667\\
Noto & Raw / LoRA / Full / Diag.
& 0.2087 / 0.3495 / 0.3112 / 0.2820
& 0.3333 / 0.3800 / 0.3400 / 0.3500
& 1.2779 / 1.0977 / 1.1467 / 1.1405
& 0.6867 / 0.8933 / 0.8767 / 0.6967\\
Turkey & Raw / LoRA / Full / Diag.
& 0.2273 / 0.2147 / 0.3025 / 0.2879
& 0.3400 / 0.3233 / 0.3067 / 0.3400
& 1.4333 / 1.1980 / 1.1676 / 1.2694
& 0.7300 / 1.0233 / 0.9133 / 0.8733\\
\bottomrule
\end{tabular}
\end{table*}

\section{Complete Rank-16 Alpha Sweep}
Tables~\ref{tab:fullalpha} and~\ref{tab:diagalpha} contain all 48 runs behind
the oracle summary. The non-monotonic response is important: increasing the
residual bound does not consistently improve performance, and the best alpha
differs by event and controller. Consequently, the maxima are used to diagnose
capacity, not as the fixed primary estimator.

\begin{table}[t]
\centering\scriptsize
\caption{Full rank-16 macro-F1 by alpha.}
\label{tab:fullalpha}
\begin{tabular}{lrrrrrr}
\toprule
Event & .5 & 1 & 2 & 3 & 5 & 10\\
\midrule
Hawaii & .2622 & .3338 & \textbf{.3869} & .2957 & .3239 & .3476\\
Libya & .3249 & .2760 & \textbf{.3852} & .3739 & .3394 & .2943\\
Noto & .2354 & \textbf{.3270} & .3095 & .3112 & .3021 & .2770\\
Turkey & .2627 & .1876 & .2527 & \textbf{.3025} & .2472 & .2578\\
\bottomrule
\end{tabular}
\end{table}

\begin{table}[t]
\centering\scriptsize
\caption{Diagonal rank-16 macro-F1 by alpha.}
\label{tab:diagalpha}
\begin{tabular}{lrrrrrr}
\toprule
Event & .5 & 1 & 2 & 3 & 5 & 10\\
\midrule
Hawaii & .2880 & .2972 & .3021 & .2918 & .3211 & \textbf{.3341}\\
Libya & .3110 & .3033 & .3041 & .2560 & \textbf{.3390} & .3278\\
Noto & .2629 & .2677 & .2278 & .2820 & \textbf{.3404} & .3310\\
Turkey & .2637 & .2698 & .2604 & \textbf{.2879} & .2843 & .2601\\
\bottomrule
\end{tabular}
\end{table}

\section{Supervised Ablations}
Table~\ref{tab:ablations} averages four events for the compact Full controller
at alpha 3 while changing one factor at a time. The layer result is especially
clear: middle+last substantially outperforms either layer alone. Rank 16 is the
best mean in the original rank sweep, while eight updates slightly outperform
four. Support48 does not improve on support24, consistent with support quality
and spatial coverage mattering more than count alone.

\begin{table}[t]
\centering\small
\caption{Four-event mean macro-F1 for one-factor supervised ablations.}
\label{tab:ablations}
\begin{tabular}{llr}
\toprule
Factor & Setting & Mean F1\\
\midrule
Support & 12 / 24 / 48 & .2538 / \textbf{.2964} / .2754\\
Layers & 14 / 27 / 14+27 & .1964 / .2213 / \textbf{.2964}\\
Rank & 5 / 8 / 16 / 32 & .2964 / .2972 / \textbf{.3136} / .2930\\
Updates & 1 / 2 / 4 / 8 & .2486 / .2219 / .2964 / \textbf{.3048}\\
Learning rate & .01 / .05 / .1 / .2 & .2659 / \textbf{.2964} / .2513 / .2466\\
\bottomrule
\end{tabular}
\end{table}

\section{Support-Seed Robustness}
Table~\ref{tab:seedmeans} summarizes the three balanced support24 seeds used by
the compact alpha-3 controller. Variation is substantial, particularly on Noto
Full and Libya Diagonal, motivating paired resampling over both support and
queries. The statistical comparisons reported in the main paper use 10,000
stratified bootstrap and paired permutation replicates.

\begin{table}[t]
\centering\small
\caption{Macro-F1 mean $\pm$ sample standard deviation over three support
seeds.}
\label{tab:seedmeans}
\begin{tabular}{lrr}
\toprule
Event & Full & Diagonal\\
\midrule
Hawaii & $.2795\pm.0317$ & $.2883\pm.0344$\\
Libya & $.3044\pm.0223$ & $.3111\pm.0530$\\
Noto & $.2611\pm.0566$ & $.1982\pm.0116$\\
Turkey & $.2317\pm.0498$ & $.2557\pm.0293$\\
\bottomrule
\end{tabular}
\end{table}

\section{Mechanism and Negative Controls}
The mean-gradient control keeps only the normalized aggregate gradient and is
therefore intrinsically rank one. Centered covariance subtracts the outer
product of the token-level mean, while random controls use orthonormal Gaussian
bases with three seeds. Their mean macro-F1 values are .2506, .2721, and .2414,
respectively, versus .2964 for the uncentered second moment. Centering yields a
similar NLL reduction (.1440 versus .1435) but lower F1, suggesting that both
the gradient mean and directional variation contribute.

The dense all-layer negative control changes 56 projections. With rank 16 it
has 14,336 coefficient scalars, but the original alpha-3/lr-.05 settings cause
loss spikes and yield event F1 values .2400, .1670, .1704, and .1935. On Hawaii,
alpha .3/lr .005 restores monotonically decreasing support loss yet yields only
.2923 query F1 and zero damaged recall. This rules out the simple explanation
that more controlled layers or more coefficient parameters necessarily help.

\section{Efficiency and Artifact Accounting}
\begin{table}[t]
\centering\small
\caption{Exact Hawaii query300 efficiency on one RTX 3090.}
\label{tab:efficiency}
\begin{tabular}{lrrr}
\toprule
Method & sec/sample & Peak MiB & Trainable\\
\midrule
Raw VLM & 1.3405 & 17,470 & 0\\
LoRA & 1.3725 & 17,508 & 10,092,544\\
Full rank5 & 1.3382 & 17,470 & 100\\
Diagonal rank5 & 1.3384 & 17,470 & 20\\
\bottomrule
\end{tabular}
\end{table}

The timing excludes adaptation and measures identical query inference. KV basis
extraction averages about 17 seconds and coefficient fitting about 53 seconds.
The compact state files are approximately 44 KiB because they contain the
bases as well as coefficients; the LoRA adapter is approximately 39.5 MiB. Full
rank-16 has 1,024 learned coefficient scalars and Diagonal rank-16 has 64, but
their exact end-to-end timing was not substituted for the measured compact
benchmark in Table~\ref{tab:efficiency}.

\section{Reproducibility and Selection Caveats}
All comparisons within a row use identical saved support and query manifests,
model fingerprints, crop size, and query IDs. Completed predictions are
configuration-bound and resumable. Nevertheless, two limitations affect
interpretation. First, the current support sets are class-balanced but spatially
concentrated (three Hawaii support tiles and four Noto support tiles). Second,
two nominally matching Hawaii rank-16 Full alpha-3 runs differ by .0338 F1,
despite matching recorded configuration. Independent repeats are therefore
required for sensitive oracle peaks. Finally, no per-event alpha selected from
the query sweep is presented as a test-independent primary result.

\bibliographystyle{IEEEtran}
\bibliography{refs}
\end{document}
